\chapter{Neutron Diffraction}
Neutron diffraction has several major advantages over x-rays, except that neutrons are much more difficult to generate and, even at the world's most powerful sources, have much lower fluxes than available routinely from x-rays. Thus, the number one goal when designing neutron instrumentation is to preserve as much of the neutron beam as can be done. This typically means that neutron beams are big -- on the scale of centimeters, though samples are sometimes sub-millimeter. Shielding for neutron instruments is also massive. Typical shielding contains hydrogen-containing materials, which moderate high-energy neutrons into thermal neutrons and neutron-absorbing materials, such as boron, cadmium, or gadolinium, which will absorb thermal neutrons effectively. Such shielding is typically massive with many tons of shielding blocks. Measurement times traditionally have been long, hours to a day, but with some modern instruments and larger samples, measurements may be completed in minutes. 

Note that neutrons readily penetrate most metals, which allows furnaces, cryostats and sample refrigerators to be readily used at neutron diffractometers. It is common to collect data at temperatures in the range of 4\ K to 10\ K as these are easy to achieve and provide for minimal loss of high Q intensity from the Debye-Waller factor, but use of neutrons can be useful for high-temperature studies as well. Many forms of operando and \textit{in situ} experiments benefit from the rapid measurements available from synchrotrons, but where possible neutrons should be considered. GSAS-II allows neutron diffraction analysis to be combined with x-ray diffraction. Even x-ray single-crystal diffraction studies can benefit from a combined refinement with a neutron powder diffraction dataset. For complex parametric studies it may well make sense to obtain neutron powder diffraction data under selected conditions. 

There are two very different types of neutron diffraction instrumentation: Those based on pulsed spallation sources where neutrons are energy-separated by the neutron's velocity. These are called time-of-flight (TOF) diffractometers and the pattern is recorded as intensity as a function of arrival time for the neutron, which can be converted to Q. Nuclear reactors create a broad spectrum neutron source where a specific wavelength is selected and used for a conventional, angular dispersive instrument where intensity is recorded as a function of angle. These are known as constant wavelength (CW) neutron diffractometers. Hybrids between the two technologies are possible. 

Neutrons are scattered from atom nuclei and the strength of this scattering is independent of Q. The scattering strength is dependent of on isotope of the element, though the cost of isotopically separated materials often prevents this from being exploited. Nonetheless, use of deuterium (${}^2\rm{H}$ is common for neutron experiments as the ${}^1\rm{H}$ component in natural abundance hydrogen contributes very large amounts of incoherent and inelastic scatter, which appears as background and can require use of a correction for absorption. Another example is that boron is a very strong absorber of neutrons and natural-abundance boron-containing samples are usually two absorbing for diffraction measurements. However, ${}^{10}\rm{B}$ is not a neutron absorber and can be used in neutron diffraction without concerns, other than for cost. 

For magnetic materials, neutrons offer one additional benefit. Neutrons are scattered from unpaired electrons, though with a scattering factor that falls-off with Q faster than the form-factor for x-ray scattering. When a material undergoes a magnetic transition (usually at lowered temperatures), additional intensity will be added to the pattern, often at new peak locations from the magentically scattering atoms. 

The processes that produce neutrons, usually create high energy neutrons (MeV) but neutrons needed for diffraction must be of much lower energy ( circa $10^{-2}$ eV or equivalently $\approx$10 meV ). The energy of neutrons is lowered by allowing them to scatter inelastically, typically from hydrogen-containing materials. This process is called moderation. These neutrons, with lowered energy, are typically called thermal neutrons, as a 25 meV neutron has a temperature of $\approx 290$ K.

The facilities that host neutron instrumentation usually require that you write a proposal to be granted time at the facility. Before submitting a proposal, it is worthwhile to contact a instrument scientist at the facility to discuss your intended project. This will help you verify that your intended measurement is feasible and that you have selected the optimum instrument for the work. 

\section{Reactor-based instrumentation}
Nuclear reactors produce neutrons from fission of ${}^{235}\rm{U}$. Fission produces neutrons with very high energies and these neutrons must be moderated at least partially to allow for the nuclear chain reaction needed to sustain the fission process. In a reactor neutron source, special attention is given to the moderator design to ensure that no higher energy (epithermal) neutrons remain in the spectrum that the reactor is designed to emit. Some neutron facilities also have a secondary moderator that is cooled to cryogenic temperatures to produce very low energy (long wavelength) ``cold'' neutrons. Other than for magnetic scattering, cold neutrons are not very useful for powder diffraction as the maximum Q value that can be obtained is too low. In fact, the ideal neutron source for reactor-based powder diffraction would have a moderator operating at a somewhat elevated temperature, should someone with very deep pockets ever wish to design an optimal instrument.

The first component in a neutron diffractometer is typically a Soller collimator. These are long thin blades arranged vertically that limit the horizontal divergence of the neutron beam, but do not limit the vertical divergence. The second component is a monochromator, but there are several differences between the monochromator used for neutrons vs. synchrotrons though commonly the same materials, such as silicon and germanium are commonly used. The first difference is that while synchrotrons typically use perfect crystals which offer very high angular and energy resolution, neutron monochromators use highly imperfect crystals that have degraded angular and energy resolution. The intensity loss from higher resolution cannot be afforded. Creating these highly imperfect crystals is something of an art. One needs to take perfect crystals and create highly mosiac domains through some form of mechanical process. Secondly, usually the monochromator is constructed from multiple crystals and these are arranged to focus the horizontally diverging beam back onto the sample. 

There is likely to be a second Soller collimator between the monochomator and the sample. Samples may be as large as finger-sized (1 cm by 8 cm) and are commonly placed in vanadium cans, though aluminum is also common. Vanadium at natural abundance has almost no coherent scattering, but does add some background to the diffraction pattern. Aluminum has significant Bragg scattering, but depending on the wavelength this will be in a fairly small number of peaks, but Al does not add any significant amount of background. 

Most instruments place a large number of detectors around the sample (more than 20 typically) with a third Soller collimator between the sample and the detector. Detection is typically done with discrete ${}^3\rm{He}$ tubes, which produce an electronic pulse similar to a scintillator detector for x-rays, but other types of neutron detectors exist, including linear position sensitive detectors and area detectors. 

The diffractometer will have a resolution curve that has optimum performance when the $2\theta$ angle for the diffractometer is about the same as the  $2\theta$ angle for the monochromator and it is common to design the instrument to have this angle in the range of 90 degrees. 

CW diffraction patterns appear not very different from x-ray diffraction patterns, though closer examination will note that the intensities fall off with $2\theta$ at a much slower rate than with x-rays due to the flat scattering factors for neutrons. Peaks are usually fairly broad with Gaussian peaks that often do not show any sample-broadening effects. One may note that the instrumental resolution will improve with Q. CW neutrons can offer quite good resolution, but only for a limited Q range. 

\section{Accelerator-based (TOF) instrumentation}

The other way that neutrons can be produced at great enough levels to allow powder diffraction is through spallation, which is the term initially used for sand-blasting, where material is machined by impacting it with small abrasive particles. In the case of neutron spallation, protons are accelerated into a target. The target is made from a high-Z metal, where the nuclei are rich with neutrons. The collision between the proton beam and the metal nucleus displaces protons from the atoms, but these neutrons are of very high energy and must be moderated before they can be used for scattering experiments. 

It is much easier to design and build an electron accelerator than a proton accelerator, and indeed early work in neutron spallation was done with electron beams, but protons are much more effective due to their greater mass. A proton accelerator could be designed to produce an almost continuous stream of protons, which would provide a nearly steady-state source of neutrons, but spallation sources are designed where large packets of protons are brought into impact the target at periodic intervals, on the order of ten times per second. Each packet of protons causes a nearly instantaneous pulse of fast neutrons when impact occurs, but most of the proton packet's energy is transmitted to the target in the form of heat, which must be removed by cooling so that the target does not melt. The generated neutrons must be moderated, by allowing them to partially equilibrate with a material (usually hydrogenous or deuterated). However, as will be made clear, one does not want to overly degrade the time width of pulse of neutrons, so moderator design is an art with compromises between efficiency, energy spectrum and the pulse duration. Also, since different types of neutron scattering instruments typically share a moderator, where each type of instrument may be optimized by different moderator characteristics, there may also be a compromise by instrument needs. Powder diffractometers rarely are equals in these design compromise discussions.  

As the fast neutron pulse hits the moderator, some will be transmitted very quickly with very little moderation. Usually a spinning drum, called a chopper, placed between the moderator and the sample will prevent these neutrons from hitting the sample. The chopper will open to allow partially moderated neutrons to escape and will close before the lowest energy neutrons can escape. The duration that the chopper is open will be short and as the neutrons travel towards the sample, they spread out in time, with the highest energy (shortest wavelength) arriving first. Thus, by synchronizing the counting electronics with the neutron pulse emission and knowing the distance between the moderator and where the detector is located, one can determine the velocity of the neutron. From the de Broglie equation, $\lambda = h/{mv}$, the neutron  wavelength is be determined. While a single neutron pulse may produce only a small number of observed diffraction events, with many neutron pulses each second, the intensities in the pattern build up with data collection time. Since the neutron pulse has finite width and the will be some uncertainty in time measurement for when a neutron is detected, this translates to a peak width. Peaks tend to be sharper on the high energy side, as few neutrons will arrive earlier than the mean arrival time, but it is the behavior of moderators that neutrons continue to leave well after the most of the pulse has left, leading to asymmetric peak shapes with a tail at longer times (shorter Q). 

Since neutron data collection time is precious, large numbers of detectors are placed around the sample. The detectors do not move, since each detector sees a spectrum of neutrons and thus a section of the full diffraction pattern. The pattern shifts to higher Q for detectors at higher Bragg angles. Also, while each detector sees the same wavelength band, the higher angle detectors offer much better resolution in terms of $\Delta\rm Q/\rm Q$ . 

One would not want to deal with, say 100 diffraction patterns from an instrument with 100 detectors, so early TOF instruments had electronics that would shift the timing for neighboring detectors so that the patterns could be superimposed for a single detector in the middle of the detector group. The detector groups were known as banks. This process was known as time-focusing, but is not actually a focusing effect, merely a form of binning. One might expect problems from time-focusing, in that detectors at different angles have differing resolution and that adding higher resolution data to lower resolution data creates, well, lower resolution data with better counting statistics, but in practice in these instruments there was a small range of angles so that the resolution was about the same and resolution degradation was imperceptible. Problems could also arise because there are corrections that are sometimes wave-length dependent corrections applied in fitting (e.g. for absorption, neutron scattering edge resonances and extinction), but since the wavelength shifts are small and the shifts are clustered around the mean, these also did not introduce any significant problems. 

Time-shifting is now done in an instrument's control software and may be done over wider angular ranges. This is not such a good idea. It is convenient to have data from few banks, but degrades the highest resolution data and makes it impossible to perform wavelength-dependent computations. It is possible to apply the absorption correction as part of the data reduction process, but this requires accurate knowledge of the composition of the sample and its density. When these are not known well, it is better to have the option to refine the absorption correction as part of fitting. The other corrections become impossible. Data should be mixed only over narrow angular ranges. 

An alternate approach is to create pseudo-CW data from TOF data, by binning in Q or $2\theta$ for relatively narrow wavelength bands. These data will appear similar to CW data, but will have the asymmetric peak shape seen with TOF. This is supported in GSAS-II. 

There are not very many TOF neutron diffractometers in the world, since there are only a few spallation neutron sources, but each one is unique, as each was designed with a different set of goals. The flight path length, meaning how far the instrument is placed from the neutron source, is perhaps the most important parameter, as that dictates the resolution of the instrument, but as the path becomes longer either the wavelength range must be decreased or the number of pulses per second accepted by the instrument must be dropped. Design variables include the type and number of detectors and how far they are placed from the sample, as well as how they are distributed with respect to Bragg and azimuthal angle.

\section{Neutron absorption}
Absorption of neutrons is not usually a major problem, but there are some elements where absorption can be problematic and others where even minor levels of an element render a measurement impossible. Likewise, while not exactly absorption, incoherent scatter of neutrons renders then unable to contribute to diffraction and can be lumped in with absorption. Both should always be considered before applying for beam time. I recommend using the web calculation tool at NIST https://www.ncnr.nist.gov/resources/activation/ to estimate sample absorption and incoherent scattering. Note that the effective $\mu$ value is the 1 over the computed ``1/e penetration depth'' for ``abs+incoh'' (absorption and incoherent scattering). For TOF experiments, consider the longest wavelength that the instrument uses (with the planned chopper settings). 

As was discussed for x-ray diffraction in more depth in section \ref{ref:muR}, with $\mu R < 1.0$, absorption can be ignored. For $1.0 < \mu R < 1.5$, experiments are feasible, if not ideal. For $1.5 < \mu R < 2.5$, experiments are slow and far from ideal, but are feasible. For $2.5 < \mu R$, you are likely wasting valuable neutrons. Unlike with x-rays, your options are more limited as to how to redesign experiments to avoid absorption/incoherent scattering. Changing composition to avoid a problematic isotope is expensive, but may be possible with a few light elements (H, B, Li) but is probably not possible further down the periodic table, but you should look up prices. Changing wavelengths is only possible to a minor extent, if at all. Changing the sample diameter is readily possible, but comes at the expense of increased data collection time, but is still preferred if it can bring $\mu R \sim 1.0$. One other option to be considered with absorbing samples is use of a annular sample. This places a ring of sample around either a hollow or minimally absorbing/scattering central object. This places much more sample into the beam than container with a radius that matches the ring thickness, but the annular container will show a similar magnitude for absorption. Note that GSAS-II does not offer corrections for an annular sample container, so that correction will need to be applied to the measured intensities and their uncertainties before the data can be input to GSAS-II. 

The neutron activation values available from the NIST web site will probably not be very accurate, without help from an experienced beamline scientist or health physicist to establish the dose, but will give an idea of which elements in your sample will become radioactive and for how long. One should assume that every sample coming out from a neutron beam is radioactive and should be handled accordingly, but longer-lived neutron activation will restrict what can be done with a material, after it has been used as a neutron sample. 

\section{Neutron detection and counting statistics}
Traditionally neutrons have been detected with ${}^3$He detectors, which produces an electrical pulse when a neutron is absorbed by an atom of  ${}^3$He, but many other types of neutron detectors have also been developed. As far as I am aware, all count individual neutrons, so that counting statistics (see section \ref{ref:CountingStatistics} are obeyed. 
